% =============================================================================
%  lec23.tex —— 第 23 讲 模 I（底本 PDF 第 175–179 页）
% =============================================================================
\lecture{23}{模 I}
\begin{center}\small\color{black!55}（底本 PDF 第 175–179 页）\end{center}

本讲我们为 $\infty$-操作子 $\OO^{\otimes}$ 定义 $\OO$-模。

\section{提纲}

\nitem\label{L23-1.1} 在前几讲中，我们构造了操作子 $\mathrm{CMod}^{\otimes}$、$\mathrm{LMod}^{\otimes}$、$\mathrm{RMod}^{\otimes}$ 与 $\mathrm{BMod}^{\otimes}$，使得对称幺半 $\infty$-范畴中相应的代数对象，是由交换/结合代数及其某几类模构成的二元组/三元组。然而，这些构造是依据 $\mathrm{Comm}^{\otimes}$ 与 $\mathrm{Assoc}^{\otimes}$ 的显式描述、以权宜的方式进行的。本讲我们对连贯 $\infty$-操作子 $\OO^{\otimes}$ 给出 $\OO$-代数的一般模理论。

\nitem\label{L23-1.2} 连贯 $\infty$-操作子的定义将在后面给出。目前我们把它当作黑箱。

设 $\OO^{\otimes}$ 是连贯 $\infty$-操作子，$p\colon\cC^{\otimes}\to\OO^{\otimes}$ 是 $\OO$-幺半 $\infty$-范畴。设 $A\in\Alg_{/\OO}(\cC)$ 是 $\OO$-代数。对每个颜色 $X\in\OO$，我们将构造一个 $\infty$-范畴 $\Mod^{\OO}_{A}(\cC)^{\otimes}_{X}$，称为 $\cC$ 中颜色为 $X$ 的 $\OO$-$A$-模 $\infty$-范畴。将有一个遗忘函子
\[
\Mod^{\OO}_{A}(\cC)^{\otimes}_{X}\to\cC^{\otimes}_{X},
\]
它把 $\OO$-$A$-模送到其\textbf{底层对象}。事实上，我们将允许 $X$ 是\textbf{多色}，即 $\OO^{\otimes}$ 中的任意对象。此外，我们将有一个\textbf{参数化版本}
\begin{equation}\label{L23-eq1.1}
\Mod^{\OO}(\cC)^{\otimes}\to\Alg_{/\OO}(\cC)\times\OO^{\otimes},
\end{equation}
使得它在 $(A,X)$ 处的纤维是 $\Mod^{\OO}_{A}(\cC)_{X}$。在良好的条件下，我们有以下结果，其细节将在下一讲解释。
\begin{itemize}[leftmargin=2.2em,itemsep=0.1em]
\item 对固定的 $A\in\Alg_{/\OO}(\cC)$，\eqref{L23-eq1.1} 的纤维
\[
\Mod^{\OO}_{A}(\cC)^{\otimes}\to\OO^{\otimes}
\]
将是一个 $\OO$-幺半 $\infty$-范畴。这推广了交换代数中的相对张量积 $-\otimes_{A}-$。
\item 投影
\[
\Mod^{\OO}(\cC)^{\otimes}\to\Alg_{/\OO}(\cC)
\]
将是笛卡尔纤维化。这推广了交换代数中的限制函子。
\item 但要注意：对 $\Alg_{/\OO}(\cC)$ 中的态射 $A\to B$，反变传输函子
\[
\res\colon\Mod^{\OO}_{B}(\cC)^{\otimes}\to\Mod^{\OO}_{A}(\cC)^{\otimes}
\]
\textbf{不是} $\OO$-幺半函子。它只是\textbf{右 lax $\OO$-幺半的}，即一个 $\infty$-操作子映射。这推广了交换代数中相应的事实。
\item 对固定的 $X\in\OO^{\otimes}$，\eqref{L23-eq1.1} 的纤维
\[
\Mod^{\OO}(\cC)^{\otimes}_{X}\to\Alg_{/\OO}(\cC)
\]
将是余笛卡尔纤维化。这推广了交换代数中的诱导函子。
\item 但要注意，函子
\[
\ind\colon\Mod^{\OO}_{A}(\cC)^{\otimes}_{X}\to\Mod^{\OO}_{B}(\cC)^{\otimes}_{X},\qquad X\in\OO^{\otimes}
\]
\textbf{不能}装配成 $\OO^{\otimes}$ 上的函子。事实上，作为右 lax $\OO$-幺半函子 $\res$ 的左伴随，$\ind$ 一般只应是左 lax $\OO$-幺半的。在交换的情形中，诱导函子恰好是 $\OO$-幺半的，但对 $\OO^{\otimes}=\mathrm{Assoc}^{\otimes}$ 这已经不成立。见下面的练习~\ref{L23-1.4}。
\end{itemize}

\begin{warning}\label{L23-1.3}
对 $\OO^{\otimes}=\mathrm{Assoc}^{\otimes}$，$\Mod^{\mathrm{Assoc}}_{A}(\cC)$ 中的对象是 $(A,A)$-双模，而不是左 $A$-模或右 $A$-模。确实，要定义 $A$ 上的相对张量积，我们需要一个左 $A$-模与一个右 $A$-模。
\end{warning}

\begin{exercise}\label{L23-1.4}
设 $A\to B$ 是结合环之间的同态。证明限制函子
\[
\res\colon {}_{B}\mathrm{BMod}_{B}\to{}_{A}\mathrm{BMod}_{A}
\]
有左伴随
\[
\ind\colon {}_{A}\mathrm{BMod}_{A}\to{}_{B}\mathrm{BMod}_{B},
\]
后者是左 lax 幺半的，但不是幺半的。
\end{exercise}

\section{连贯 $\infty$-操作子}

\nitem\label{L23-2.1} $\infty$-范畴 $\Mod^{\OO}(\cC)^{\otimes}_{X}$ 将定义为某个全子 $\infty$-范畴
\[
\Mod^{\OO}(\cC)^{\otimes}_{X}\subseteq\Fun_{\OO^{\otimes}}(K_{X},\cC^{\otimes}),
\]
其中 $K_{X}$ 是 $\OO^{\otimes}$ 上的某个 $\infty$-范畴，在 Comm 情形中扮演 $\mathrm{CMod}^{\otimes}$ 的角色。

我们将通过 $K_{X}$ 的函子性得到 $\Mod^{\OO}(\cC)^{\otimes}_{X}$ 关于 $X$ 的函子性。换言之，我们将定义一个带有两个投影
\[
(\mathrm{ev}_{0},\mathrm{ev}_{1})\colon K_{\OO}\to\OO^{\otimes}\times\OO^{\otimes}
\]
的 $\infty$-范畴 $K_{\OO}$，并定义
\[
K_{X}\defeq\{X\}\times_{\OO^{\otimes},\mathrm{ev}_{0}}K_{\OO}\xrightarrow{\ \mathrm{ev}_{1}\ }\OO^{\otimes},
\]
其中纤维积使用 $\mathrm{ev}_{0}$，而投影 $K_{X}\to\OO^{\otimes}$ 使用 $\mathrm{ev}_{1}$。

为实现这一点，我们不会把 $K_{X}$ 做成 $\infty$-操作子。当 $\OO^{\otimes}=\mathrm{Comm}^{\otimes}$ 且 $X=\langle1\rangle$ 时，这意味着我们只允许 $\mathrm{CMod}^{\otimes}$ 中至多含一个 $m$-项的对象。

\nitem\label{L23-2.2} 为定义 $K_{\OO}$，我们需要一些术语。

\begin{definition}\label{L23-2.3}
设 $\OO^{\otimes}$ 是 $\infty$-操作子。设态射 $f\colon(X_{1},\dots,X_{m})\to(Y_{1},\dots,Y_{n})$ 位于 $\langle m\rangle\xrightarrow{\alpha}\langle n\rangle$ 之上。若对每个 $j\in\langle n\rangle^{\circ}$，下述条件恰有一个成立，则称 $f$ 为\textbf{半惰性态射}：
\begin{enumerate}[label=(\alph*),leftmargin=2.4em,itemsep=0.1em]
\item 集合 $\alpha^{-1}(\{j\})$ 为空。
\item 集合 $\alpha^{-1}(\{j\})$ 恰含一个元素 $i$，且相应的态射 $X_{i}\to Y_{j}$ 是同构。
\end{enumerate}
\end{definition}

\begin{exercise}\label{L23-2.4}
证明：$f$ 是惰性态射，当且仅当对每个 $j\in\langle n\rangle^{\circ}$ 都有 (b) 成立。
\end{exercise}

\begin{definition}\label{L23-2.5}
设
\[
K_{\OO}\subseteq\Fun([1],\OO^{\otimes})
\]
是 $\OO^{\otimes}$ 中由半惰性态射构成的全子 $\infty$-范畴，并设
\[
K_{X}\defeq\{X\}\times_{\OO^{\otimes},\mathrm{ev}_{0}}K_{\OO}
\]
是从 $X$ 出发的半惰性态射构成的 $\infty$-范畴。
\end{definition}

\begin{definition}\label{L23-2.6}
若函子 $\mathrm{ev}_{0}$ 与 $\mathrm{ev}_{1}$ 都把 $K_{\OO}$ 中的态射送到 $\OO^{\otimes}$ 中的惰性态射，则称该态射是\textbf{惰性的}。
\end{definition}

\begin{exercise}\label{L23-2.7}
设 $\OO^{\otimes}=\mathrm{Comm}^{\otimes}$ 且 $X=\langle1\rangle$。描述范畴 $K_{X}$，并构造一个定义在 $\mathrm{Comm}^{\otimes}$ 上的典范函子 $K_{X}\to\mathrm{CMod}^{\otimes}$。证明
\[
\Alg_{\mathrm{CMod}}(\cC)\to\Fun_{\mathrm{Comm}^{\otimes}}(\mathrm{CMod}^{\otimes},\cC^{\otimes})\to\Fun_{\mathrm{Comm}^{\otimes}}(K_{X},\cC^{\otimes})
\]
是全忠实的，且其本质像由那些保持惰性态射的函子 $K_{X}\to\cC^{\otimes}$ 组成。
\end{exercise}

\nitem\label{L23-2.8} 现在可以给出连贯 $\infty$-操作子的定义了。

\begin{definition}\label{L23-2.9}
设 $\OO^{\otimes}$ 是小 $\infty$-操作子。若它满足以下条件，则称 $\OO^{\otimes}$ 是\textbf{连贯的}：
\begin{enumerate}[label=(\roman*),leftmargin=2.4em,itemsep=0.1em]
\item $\infty$-操作子 $\OO^{\otimes}$ 是带单位的。
\item 颜色的 $\infty$-范畴 $\OO$ 是空间。
\item 沿 $\mathrm{ev}_{0}\colon K_{\OO}\to\OO^{\otimes}$ 的基变换函子
\[
\mathrm{ev}_{0}^{*}\colon(\Cat_{\infty})_{/\OO^{\otimes}}\to(\Cat_{\infty})_{/K_{\OO}},\qquad
D\mapsto D\times_{\OO^{\otimes},\mathrm{ev}_{0}}K_{\OO}
\]
有右伴随 $\mathrm{ev}_{0,*}$。
\end{enumerate}
\end{definition}

\begin{proposition}\label{L23-2.10}（HA.3.3.1.12, 4.1.1.20）
$\infty$-操作子 $\mathrm{Comm}^{\otimes}$ 与 $\mathrm{Assoc}^{\otimes}$ 是连贯的。
\end{proposition}

\nitem\label{L23-2.11} 尽管我们尚未定义 $E_{k}$，这里先记下以下结果。

\begin{theorem}\label{L23-2.12}（HA.5.1.1.1）
对任意 $k\ge0$，$\infty$-操作子 $E_{k}^{\otimes}$ 是连贯的。
\end{theorem}

\begin{remark}\label{L23-2.13}
利用 HA.B.3 中的结果，可以证明 (iii) 等价于
\par\noindent\hspace*{2.4em}(iii$'$) 函子 $\mathrm{ev}_{0}\colon K_{\OO}\to\OO^{\otimes}$ 是平坦范畴纤维化。
\par
因此由 HA.3.3.2.2，我们对连贯 $\infty$-操作子的定义与 Lurie 在 HA.3.3.1.1 中的定义一致。事实上，在这些等价条件下，伴随
\[
\mathrm{ev}_{0}^{*}\colon(\Cat_{\infty})_{/\OO^{\otimes}}\rightleftarrows(\Cat_{\infty})_{/K_{\OO}}:\!\mathrm{ev}_{0,*}
\]
可以实现为（切片）Joyal 模型范畴之间的 Quillen 伴随。由此可知，对 $K_{\OO}$ 上的范畴纤维化 $D$，函子 $\mathrm{ev}_{0,*}$ 把它送到 $\OO^{\otimes}$ 上的拟范畴 $\mathrm{ev}_{0,*}(D)$，后者满足以下泛性质
\[
\Hom_{(\Set_{\Delta})_{/\OO^{\otimes}}}(J,\mathrm{ev}_{0,*}(D))
\simeq
\Hom_{(\Set_{\Delta})_{/K_{\OO}}}(J\times_{\OO^{\otimes}}K_{\OO},D),
\]
其中右端的纤维积是朴素纤维积。
\end{remark}

\begin{exercise}\label{L23-2.14}
设 $D\in(\Cat_{\infty})_{/\OO^{\otimes}}$。证明 $\mathrm{ev}_{0,*}\circ\mathrm{ev}_{1}^{*}(D)$ 在 $X\in\OO^{\otimes}$ 处的纤维等价于 $\Fun_{\OO^{\otimes}}(K_{X},D)$。
\end{exercise}

\section{模的定义}

\nitem\label{L23-3.1} 我们终于可以构造 $\Mod^{\OO}(\cC)^{\otimes}$ 了。

\begin{definition}\label{L23-3.2}
设 $\OO^{\otimes}$ 是连贯 $\infty$-操作子，$\cC^{\otimes}\to\OO^{\otimes}$ 是 $\infty$-操作子映射。我们把
\[
\Mod^{\OO}(\cC)^{\otimes}\subseteq\mathrm{ev}_{0,*}\circ\mathrm{ev}_{1}^{*}(\cC^{\otimes})
\]
定义为由满足下述条件的对象构成的全子 $\infty$-范畴：相应的函子 $K_{X}\to\cC^{\otimes}$（见练习~\ref{L23-2.14}）保持惰性态射。
\end{definition}

\nitem\label{L23-3.3} 注意 \eqref{L23-eq1.1} 中的一个投影是自动的：我们有一个结构函子
\[
\Mod^{\OO}(\cC)^{\otimes}\to\OO^{\otimes}.
\]

\begin{exercise}\label{L23-3.4}
设 $\OO^{\otimes}=\mathrm{Comm}^{\otimes}$。把上述函子在 $\langle1\rangle$ 处的纤维等同于 $\Alg_{\mathrm{CMod}}(\cC)$。
\end{exercise}

\nitem\label{L23-3.5} 为得到另一个投影
\[
\Mod^{\OO}(\cC)^{\otimes}\to\Alg_{\OO}(\cC),
\]
我们需要更多定义。

\begin{definition}\label{L23-3.6}
设 $\OO^{\otimes}$ 是 $\infty$-操作子。设态射 $f\colon(X_{1},\dots,X_{m})\to(Y_{1},\dots,Y_{n})$ 位于 $\langle m\rangle\xrightarrow{\alpha}\langle n\rangle$ 之上。若对每个 $j\in\langle n\rangle^{\circ}$，
\begin{enumerate}[label=(\alph*),leftmargin=2.4em,itemsep=0.1em]
\item 集合 $\alpha^{-1}(\{j\})$ 为空，
\end{enumerate}
则称 $f$ 为\textbf{零态射}。设 $K^{0}_{\OO}\subseteq K_{\OO}$ 是由零态射构成的全子 $\infty$-范畴。
\end{definition}

\yizhu{这里的“零态射”（null morphism）指可通过 $\langle0\rangle$ 分解的态射（见练习~\ref{L23-3.7}），与通常范畴论中经零对象分解的“零态射”不同；为与原文一致，仍译作“零态射”。\textbf{[来源：I]}}

\begin{exercise}\label{L23-3.7}
证明：$f$ 是零态射，当且仅当它通过对象 $\langle0\rangle\in\OO^{\otimes}$ 分解。
\end{exercise}

\begin{exercise}\label{L23-3.8}
证明函子 $(\mathrm{ev}_{0}^{0},\mathrm{ev}_{1}^{0})\colon K^{0}_{\OO}\to\OO^{\otimes}\times\OO^{\otimes}$ 是等价。
\end{exercise}

\begin{exercise}\label{L23-3.9}
证明下图
\begin{center}
\begin{tikzcd}[column sep=3.4em,row sep=2.6em]
& K^{0}_{\OO} \arrow[dl, "\mathrm{ev}_{0}^{0}"'] \arrow[d] \arrow[dr, "\mathrm{ev}_{1}^{0}"] & \\
\OO^{\otimes} & K_{\OO} \arrow[l, "\mathrm{ev}_{0}"'] \arrow[r, "\mathrm{ev}_{1}"] & \OO^{\otimes}
\end{tikzcd}
\end{center}
诱导一个自然函子
\[
\mathrm{ev}_{0,*}\circ\mathrm{ev}_{1}^{*}(\cC^{\otimes})\to
\mathrm{ev}_{0,*}^{0}\circ\mathrm{ev}_{1}^{0,*}(\cC^{\otimes})
\simeq\OO^{\otimes}\times\Fun_{\OO^{\otimes}}(\OO^{\otimes},\cC^{\otimes}),
\]
它限制为一个定义在 $\OO^{\otimes}$ 上的函子
\[
\Mod^{\OO}(\cC)^{\otimes}\to\OO^{\otimes}\times\Alg_{\OO}(\cC).
\]
\end{exercise}

\nitem\label{L23-3.10} 最后我们构造遗忘函子
\[
\Mod^{\OO}(\cC)^{\otimes}\to\cC^{\otimes}.
\]

\begin{exercise}\label{L23-3.11}
证明 $\OO^{\otimes}$ 中的任何同构都是半惰性的。特别地，存在函子
\[
\Delta\colon\OO^{\otimes}\to K_{\OO},\qquad X\mapsto\id_{X}.
\]
\end{exercise}

\begin{exercise}\label{L23-3.12}
证明下图
\begin{center}
\begin{tikzcd}[column sep=3.4em,row sep=2.6em]
& \OO^{\otimes} \arrow[d] \arrow[dl, double, no head] \arrow[dr, double, no head] & \\
\OO^{\otimes} & K_{\OO} \arrow[l, "\mathrm{ev}_{0}"'] \arrow[r, "\mathrm{ev}_{1}"] & \OO^{\otimes}
\end{tikzcd}
\end{center}
诱导一个定义在 $\OO^{\otimes}$ 上的自然函子
\[
\mathrm{ev}_{0,*}\circ\mathrm{ev}_{1}^{*}(\cC^{\otimes})\to\cC^{\otimes}.
\]
其限制
\[
\Mod^{\OO}(\cC)^{\otimes}\to\cC^{\otimes}
\]
称为\textbf{遗忘函子}。
\end{exercise}

\appendixsection{平坦纤维化}

\nitem\label{L23-A.1} 设 $q\colon X\to Y$ 是拟范畴之间的范畴纤维化。正如注~\ref{L23-2.13} 中提到的，函子 $q_{*}\colon(\Cat_{\infty})_{/X}\to(\Cat_{\infty})_{/Y}$ 存在，当且仅当 $q$ 是平坦范畴纤维化。后一条件通过单纯集有一个很好的刻画：$q$ 是平坦的，当且仅当对任意 $\Delta^{2}\to Y$，映射
\[
X\times_{Y}\Lambda^{2}_{1}\to X\times_{Y}\Delta^{2}
\]
是范畴等价。非正式地说，这意味着对 $Y$ 中的任意链 $u\xrightarrow{f}v\xrightarrow{g}w$，$X_{u}$ 与 $X_{w}$ 之间在 $g\circ f$ 上的对应，是 $f$ 上的对应与 $g$ 上的对应的复合。

\nitem\label{L23-A.2} \textbf{建议阅读}\quad HA.B.3。
